% Abstract — write last. Target 200-250 words for PLOS Comp Bio,
% up to 350 for Frontiers in Comp Neuro.

\begin{abstract}
\noindent
\todo{Write last. Suggested structure:}
\begin{itemize}
  \item[(1)] \textbf{Context.} Computational spinal CPG models in the
        Rybak/Danner/Zhang lineage capture central interlimb coordination
        in fine detail but, by the authors' own account, omit motoneurons,
        muscle dynamics, sensory feedback, and synaptic plasticity.
  \item[(2)] \textbf{Question.} Can a computational sCPG self-organise its
        descending and cutaneous afferent weights through STDP while
        being closed-loop with biomechanical output?
  \item[(3)] \textbf{Approach.} Two-leg NEST model of a rat spinal CPG
        with Izhikevich neurons, STDP on BS$\to$RG and CUT$\to$RG synapses,
        closed-loop Ia$\to$reciprocal-inhibition feedback, motoneuron pools,
        Hill-like force/length proxies, and externally paced heel$\to$toe
        Ia activation.
  \item[(4)] \textbf{Results.} Stable counter-phase locomotion for
        \SI{120}{\second} (120 cycles), peak force 17.4\,a.u., E/F
        correlation $-0.91$. Cycle period tracks commanded step period
        across \SIrange{600}{1400}{\milli\second} (rat walk$\to$trot range)
        with $<\SI{0.5}{\milli\second}$ error. Ablation reveals that
        commissural inhibition and paced descending drive are necessary;
        the Ia loop and asymmetric reciprocal inhibition are redundant
        under paced drive.
  \item[(5)] \textbf{Implications.} Locomotor patterns can self-organise
        from STDP given paced sensory drive, extending the Rybak central
        architecture below the rhythm generator and into the closed loop.
\end{itemize}
\end{abstract}

\noindent\textbf{Keywords:} spinal central pattern generator, STDP,
locomotion, sensory feedback, Ia afferent, NEST simulator, computational
neuroscience.
